% ═══════════════════════════════════════════════════════════════════════════════
%  CAPÍTULO — Facade
% ═══════════════════════════════════════════════════════════════════════════════

\chapter{Facade}

\patroninfo{Estructural}{185}

% ─── Situación Problema ──────────────────────────────────────────────────────
\section{Situación Problema}

El proceso de confirmar una reserva en Airbnb involucra internamente
varios subsistemas independientes: el \emph{motor de disponibilidad}
verifica que las fechas estén libres, el \emph{servicio de pagos}
autoriza y retiene el importe, el \emph{servicio de notificaciones}
avisa al anfitrión, el \emph{generador de contratos} produce el acuerdo
de estancia y el \emph{sistema de calendario} bloquea las fechas.

Cada subsistema expone su propia API con distintas convenciones.
El código del cliente —la pantalla de confirmación de la aplicación móvil—
debe conocer y orquestar todos estos sistemas en el orden correcto, lo
que genera un acoplamiento elevado y hace que cualquier cambio en un
subsistema repercuta directamente en la capa de presentación.

% ─── Solución ────────────────────────────────────────────────────────────────
\section{Solución}

El patrón Facade introduce la clase \texttt{ServicioReserva}, que ofrece
un único método \texttt{confirmarReserva()} de alto nivel. Internamente,
esta fachada coordina la secuencia correcta de llamadas a los subsistemas:
verifica disponibilidad, autoriza el pago, genera el contrato, bloquea
las fechas y dispara las notificaciones.

La capa de presentación solo conoce \texttt{ServicioReserva} y queda
desacoplada de los detalles internos. Si el equipo de backend modifica
el orden de los pasos o incorpora un nuevo subsistema —como un servicio
de detección de fraude— solo es necesario actualizar la fachada, sin
tocar la interfaz que consume el cliente.

% ─── Diagrama de clases ─────────────────────────────────────────────────────
\begin{figure}[H]
    \centering
    % \includegraphics[width=\textwidth]{imagenes/abstract_factory_diagrama.png}
    \caption{Diagrama de clases del patrón Abstract Factory}
\end{figure}


% ─── Roles de las Clases ────────────────────────────────────────────────────
\section{Roles de las Clases}

% Explicar la responsabilidad de cada clase/interfaz

\begin{itemize}
    \item \textbf{AbstractFactory:}
    \item \textbf{ConcreteFactory:}
    \item \textbf{AbstractProduct:}
    \item \textbf{ConcreteProduct:}
    \item \textbf{Client:}
\end{itemize}


% ─── Main ───────────────────────────────────────────────────────────────────
\section{Main}

% Explicar qué realiza el programa principal

\begin{lstlisting}[language=Java, caption=NombreClase]
 
\end{lstlisting}


% ─── Salida del Main ────────────────────────────────────────────────────────
\section{Salida del Main}

\begin{lstlisting}[language=Java, caption=NombreClase]
 
\end{lstlisting}


% ─── Clases ─────────────────────────────────────────────────────────────────
\section{Clases}

% ─── Clase 1 ────────────────────────────────────────────────────────────────
\subsection{NombreClase}

\begin{lstlisting}[language=Java, caption=NombreClase]
 
\end{lstlisting}


% ─── Clase 2 ────────────────────────────────────────────────────────────────
\subsection{NombreClase}

\begin{lstlisting}[language=Java, caption=NombreClase]
 
\end{lstlisting}


% ─── Clase 3 ────────────────────────────────────────────────────────────────
\subsection{NombreClase}

\begin{lstlisting}[language=Java, caption=NombreClase]
 
\end{lstlisting}


% ─── Conclusiones ───────────────────────────────────────────────────────────
\section{Conclusiones}

% Explicar:
% - Ventajas del patrón
% - Cuándo usarlo
% - Beneficios obtenidos
% - Posibles desventajas
% - Relación con principios SOLID